% Main-paper replacement for the supplied Effect System subsection.
% Companion: effect-system-appendix.tex, included after \appendix.
% These are manuscript fragments; use the paper's existing bibliography.

\subsection{Effect System}
\label{sec:effect-system}
Our effect system identifies spatial and temporal parallelism in graph
loops and guides their realization. Its central abstraction records both
the locations accessed by a computation and its phase and occurrence
scope within a traversal. This structure lets the compiler reason about
ownership of graph work, ordering between phases, and dependences across
rounds using the same effect representation.

\paragraph{Phase- and scope-aware effects.}
A canonical graph loop traverses a frontier or the full vertex domain
$F_t$. For each source $u$, it executes a preamble $P(u)$, a body $B(u,v)$
for each neighbor $v\in N(u)$, and an epilogue $Q(u)$:
\begin{equation}
E_t=\mathrm{Seq}_{u\in F_t}\Bigl(
P(u)\triangleright
\bigl(\mathrm{Seq}_{v\in N(u)}B(u,v)\bigr)
\triangleright Q(u)\Bigr).
\label{eq:effect-expr}
\end{equation}
Here $\mathrm{Seq}$ retains serial enumeration order and
$\triangleright$ records phase order within a source iteration.
Within each phase, effects form a multiset, grouped by structural
composition $\parallel$; this operator does not assert concurrency.
Our use of sequential composition follows sequential effect
systems~\cite{10.1145/2429069.2429074} and is inspired by effect
quantales~\cite{gordon2017genericapproachflowsensitivepolymorphic,10.1145/3450272}.

An atom records an operation kind, canonical memory base, index region,
and occurrence scope. Kinds distinguish reads, writes, old-value-dependent
updates, recognized algebraic updates, claims, and frontier activations.
Regions distinguish source-derived ($U$), destination-derived ($V$),
data-derived ($D$), and global ($G$) accesses; scopes distinguish
once-only, per-source, and per-pair operations. Fixed-point provenance
analysis over SSA values supplies the regions. Region tags alone do not
prove that accesses are disjoint. Unsupported operations prevent the
parallel rewrite.

The compiler derives access footprints and ordered RAW, WAR, and WAW
relations, including conflicts on shared runtime resources. Temporal
classification uses endpoint provenance, the mutation's phase, and
whether traversal uses a frontier. Per-source phase order alone cannot
separate one source's preamble read from another source's neighbor-body
write. Likewise, effects grouped within a phase remain potentially
conflicting unless further analysis establishes separation or ordering.
Appendix~\ref{app:effects-extraction} gives the extraction and
classification rules.

\paragraph{Independence and correctness.}
Spatial analysis considers work instances and proposed ownership
partitions; temporal analysis considers rounds or components of source
phases. For either axis $A\in\{S,T\}$, let $X$ be the analyzed execution
units, $\mathsf{Dep}_A$ their remaining dependence relation, and
$\mathcal P$ a proposed non-trivial partition of $X$. The independence
criterion is
\begin{equation}
\begin{gathered}
\mathsf{Dep}_A\cap\mathsf{Cross}(\mathcal P)=\varnothing,\\
\mathsf{Cross}(\mathcal P)=
\bigcup_{\substack{P_i,P_j\in\mathcal P\\i\ne j}}P_i\times P_j.
\end{gathered}
\label{eq:effects-independence}
\end{equation}
Thus, no remaining dependence may cross concurrently executed parts.
For example, source ownership keeps a source's $A[u]$ updates together,
but updates to $A[v]$ can still conflict across source owners and require
additional justification.

Applicable transformations, including snapshotting, staged claims, and
privatized updates or reductions, can discharge conflicts. Each requires
both preservation of the original observable behavior and independence
of the resulting work. A snapshot must provide the value required by
the original computation; recognizing an algebraic update does not by
itself justify regrouping it. Privatization additionally requires a valid
combine operation and no unsupported observation of uncombined state.
The compiler removes a conflict only after establishing the relevant
transformation conditions.

Disconnected components of the remaining interference graph identify
independent groups. With conservative footprints and behavior-preserving
transformations, operations in different groups commute because neither
writes a location read or written by the other. A realization can
therefore execute those groups concurrently while preserving each
group's required order, occurrence counts, and control and lifecycle
constraints. Independence on one axis permits concurrency along that
axis while retaining constraints on the other.
Appendix~\ref{app:effects-independence} states the assumptions and
expands this conditional preservation argument.

\paragraph{From effects to execution.}
Given loop $L$, effects $E_t$, and analysis results and justified
transformations $\mathcal A$, the compiler constructs an execution plan
$\mathcal R(L,E_t,\mathcal A)=\Pi_t$. The plan specifies work units,
their domains and state, and precedence constraints. If no supported
realization satisfies these requirements, execution remains serial.

Spatial realization uses Graptor's CleanCut
layout~\cite{10.1145/3392717.3392753} to partition work by source or
destination ownership. Source ownership keeps a source's phases
together. Destination ownership can split its neighbors across
partitions; supported source-coverage passes preserve per-source work,
including for sources with no neighbors. Temporal realization executes
independent phase components within a source iteration or independent
groups of rounds, preserving order within each component or group.

When round independence and spatial partitioning both hold, the compiler
combines them. Each round $t$ has setup $S_t$, partition batches
$B_{t,j}$, and finalization $F_t$, ordered as
\[
S_t\longrightarrow B_{t,j}\longrightarrow F_t.
\]
Independent rounds use separate mutable auxiliary state and may overlap.
The implemented combination supports round and spatial parallelism;
source-phase parallelism uses its separate realization.

The runtime schedules executable tasks under these precedence
constraints. A task may contain many effect occurrences: the execution
DAG is distinct from the interference graph. CPU and GPU backends use
the same analysis results, with additional device-support checks for
GPU execution. Currently, each GPU temporal component executes
sequentially in one thread, exposing temporal concurrency between
kernels. Unsupported offload uses a supported CPU realization.
Appendix~\ref{app:effects-runtime} details callbacks, synchronization,
resource lifetimes, and backend execution.
